\section{A supply of nonsparse additive transforms}\label{sec:supply}

We retain the sets \(S_p\), their densities \(\sigma_p\), and the normalized
functions \(F_p,g_p\) defined in Corollary~\ref{cor:mixed-flatness}.
The normalizations give
\[
 \sum_{x\in\F_p}|F_p(x)|^2=p,\qquad
 \frac1p\sum_{v\in\F_p}|g_p(v)|^2=1,\qquad
 g_p(0)=0,\qquad g_p(-v)=\overline{g_p(v)}.
\]
The conclusion below is a separate consequence of the sieve estimates and
prime coverage. It does not require the decorrelation assertion in
\eqref{eq:src27}.
Write
\[
 \gamma_p=\frac1p\sum_{v\in\F_p}|g_p(v)|
\]
for the probability $L^1$ norm of the normalized additive transform.

\begin{proposition}\label{prop:supply}
There is an absolute constant \(\delta_0>0\) such that, for all sufficiently
large \(L\),
\begin{equation}\label{eq:src28}
 \sum_{\substack{.05L\le \log\log p\le .9L\\
                  1/3\le \sigma_p\le2/3\\
                  \gamma_p\ge\delta_0}}
        \frac1p \ \ge\ .15L .
\end{equation}
\end{proposition}

Suppose, to the contrary, that primes with small $\gamma_p$ have large
harmonic mass. We shall construct a nonnegative weight from their sparse
Fourier spectra. A tensor estimate bounds its sum over pairs of summand
elements, while a prime-distribution estimate evaluates its sum over
large primes. Prime coverage forces a lower bound for the former sum,
and the contraction furnished by the sparse spectra will make the bounds
incompatible.

\subsection{A budget for centered tensor coordinates}

Throughout this section we reset the large parameters by
\[
 \log X=\exp L,\qquad R=\sqrt X.
\]
Set
\[
 A_0=A\cap(R+N_*,X],\qquad
 D_0=D\cap[-X,-R-N_*),\qquad m_A=|A_0|,\quad m_D=|D_0|.
\]
Both sets are nonempty for large \(L\), by \eqref{eq:src1}. For every prime
\(p\le R\), their reductions are supported on \(S_p\) and
\(T_p:=S_p^c\), respectively. Write
\(\kappa_p=|T_p|/|S_p|=(1-\sigma_p)/\sigma_p\).

Apply Lemma~\ref{lem:collision} at
\(Y=X(\log X)^{40}\), using the uniform measures on its two full tails.
By \eqref{eq:src1}, their largest point masses are at most
\((\log Y)^4/\sqrt Y\) for large \(Y\). The associated prime cutoff is
\(\sqrt Y/(\log Y)^5>R\). Since
\(\sigma_p^{-1}+(1-\sigma_p)^{-1}-4
 =\kappa_p+\kappa_p^{-1}-2\), it follows that
\begin{equation}\label{eq:src29}
 \sum_{p\le R}(\kappa_p+\kappa_p^{-1}-2)\frac{\log p}{p}\ll L,
 \qquad
 \sum_{p\le R}\frac{|\log\kappa_p|}{p}\ll\sqrt L .
\end{equation}
For the second assertion, use
\((\log t)^2\le t+t^{-1}-2\) for \(t>0\) and Cauchy--Schwarz:
\[
 \sum_{p\le R}\frac{|\log\kappa_p|}{p}
 \le
 \left(\sum_{p\le R}
       \frac{(\log\kappa_p)^2\log p}{p}\right)^{1/2}
 \left(\sum_p\frac1{p\log p}\right)^{1/2}.
\]
The last prime sum converges.

Let \(\mathbf e_x\) be the standard point vector in \(\C^{\F_p}\).
Let \(Q_{S_p}\) be the orthogonal projection onto
\[
 \mathcal H_{S_p}
 =\left\{f\in\C^{\F_p}:
      \supp f\subseteq S_p,\ \sum_x f(x)=0\right\},
\]
and define \(Q_{T_p}\) and \(\mathcal H_{T_p}\) similarly.
All these inner products use counting measure. For squarefree \(t\)
whose prime factors are at most \(R\), put
\[
 B_A(t)=
 \left\|\frac1{m_A}\sum_{a\in A_0}
       \bigotimes_{p\mid t}Q_{S_p}\mathbf e_{a\bmod p}\right\|^2,
 \qquad B_A(1)=1,
\]
and define \(B_D(t)\) with \(D_0,T_p\).

The precise energy identity we need is
\begin{equation}\label{eq:supply-primitive-energy}
 \sum_{\substack{h\bmod q\\(h,q)=1}}
       \left|\frac1{m_A}\sum_{a\in A_0}e_q(ha)\right|^2
   =\sum_{t\mid q}tB_A(t)\prod_{p\mid q/t}\kappa_p
       \qquad(q\le R,\ q\ {\rm squarefree}).
\end{equation}
Indeed, the Ramanujan kernel factors over the primes of \(q\).
For \(x,y\in S_p\),
\[
 \sum_{h\in\F_p^*}e_p(h(x-y))
 =p\ind_{x=y}-1
 =\kappa_p+
        p\langle Q_{S_p}\mathbf e_x,Q_{S_p}\mathbf e_y\rangle .
\]
Expand the product of these identities and average over two independent
uniform elements of \(A_0\). The term with centered factors at precisely
the primes of \(t\) is the corresponding term on the right of
\eqref{eq:supply-primitive-energy}. For \(D_0\), the same identity holds
with \(\kappa_p^{-1}\). In particular all terms in both expansions are
nonnegative.

Write
\[
 \mathcal B_L=\{p:.05L\le\log\log p\le.9L\}.
\]
For each fixed \(C_*>0\), let \(\mathcal T_{C_*}\) be the set of products
of at most \(C_*L\) distinct primes from \(\mathcal B_L\), including 1.
We claim that
\begin{equation}\label{eq:src30}
 \sum_{t\in\mathcal T_{C_*}}B_A(t)
       \le \exp(o(L))\frac R{m_A},
 \qquad
 \sum_{t\in\mathcal T_{C_*}}B_D(t)
       \le \exp(o(L))\frac R{m_D}.
\end{equation}
The error terms are uniform over \(t\) for each fixed \(C_*\).
To prove this, note first that
\[
 \log t\le C_*L\exp(.9L)=o(\log X).
\]
Consequently \(t<R\), and \(U=R/t=X^{1/2-o(1)}\), uniformly.
Among the positive integers \(u\le U\), the proportion failing to be
squarefree is at most \(\sum_p p^{-2}<1\). The proportion not coprime to
\(t\) is at most
\[
 \sum_{p\mid t}\frac1p
 \le C_*L\exp\bigl(-\exp(.05L)\bigr)=o(1).
\]
Also, by \eqref{eq:src29},
\[
 \frac1{\lfloor U\rfloor}\sum_{u\le U}
           \sum_{p\mid u}|\log\kappa_p|
 \ll\sum_{p\le U}\frac{|\log\kappa_p|}{p}
 \ll\sqrt L .
\]
Markov's inequality shows that only \(O(L^{-1/4})\) of these integers
have the inner sum exceeding \(L^{3/4}\). Thus at least \(cU\)
integers, for an absolute \(c>0\) and all sufficiently large \(L\), are
squarefree, coprime to \(t\), and satisfy
\[
 \prod_{p\mid u}\kappa_p\ge e^{-L^{3/4}},
 \qquad
 \prod_{p\mid u}\kappa_p^{-1}\ge e^{-L^{3/4}}.
\]
The additive large sieve, applied to the probability measure on \(A_0\),
bounds the sum of the left side of
\eqref{eq:supply-primitive-energy} over squarefree \(q\le R\) by
\(CR^2/m_A\). Interchanging the nonnegative terms on the right, and
retaining only \(t\in\mathcal T_{C_*}\), gives
\[
 \frac{CR^2}{m_A}
 \ge
 \sum_{t\in\mathcal T_{C_*}}tB_A(t)
     \sum_{\substack{u\le R/t\\u\ {\rm squarefree}\\(u,t)=1}}
          \prod_{p\mid u}\kappa_p
 \ge cR e^{-L^{3/4}}
          \sum_{t\in\mathcal T_{C_*}}B_A(t).
\]
This proves the first assertion, and the reciprocal argument proves
the second. Notice in particular the consequence of the \(t=1\) terms:
\begin{equation}\label{eq:supply-tail-size}
 m_A,m_D\le R\exp(o(L)).
\end{equation}

\subsection{A contracting local kernel from a sparse transform}

Fix a sufficiently small absolute \(\eps_0>0\); all conditions on its
size below are absolute. Choose \(0<\delta_0\le\eps_0^2\).
Suppose, for a sequence of arbitrarily large \(L\), that
\eqref{eq:src28} fails. Let
\[
 \mathcal S=\{p\in\mathcal B_L:
                  1/3\le\sigma_p\le2/3,\ \gamma_p<\delta_0\},
 \qquad H=\sum_{p\in\mathcal S}\frac1p.
\]
Mertens' theorem gives \(\sum_{p\in\mathcal B_L}1/p=.85L+o(1)\).
Outside the balanced range,
\(\kappa_p+\kappa_p^{-1}-2\ge1/2\), so \eqref{eq:src29} gives
\[
 \sum_{\substack{p\in\mathcal B_L\\
                   \sigma_p\notin[1/3,2/3]}}\frac1p
 \ll L\exp(-.05L)=o(1).
\]
It follows that \(H\ge .70L-o(L)\). We shall use the weaker bound
\begin{equation}\label{eq:supply-sparse-mass}
 H\ge .68L,
\end{equation}
which remains valid after deleting any one prime of \(\mathcal S\).

For \(p\in\mathcal S\) define
\[
 E_p=\{v\in\F_p:|g_p(v)|>1\}.
\]
This set is symmetric, avoids zero, and satisfies
\begin{equation}\label{eq:supply-concentrated-spectrum}
 |E_p|\le\eps_0p,\qquad
 \sum_{v\in E_p}|g_p(v)|^2\ge(1-\eps_0^2)p.
\end{equation}
Indeed, \(|E_p|\le\sum_v|g_p(v)|<\delta_0p\); on its complement,
\(|g_p(v)|^2\le|g_p(v)|\). Put \(t_0=.85\) and
\[
 b_p(x)=\frac{t_0}{p}\sum_{h\in E_p}e_p(hx).
\]
These functions are real. Let \(\Pi_p\) be the orthogonal projection
in \(\C^{\F_p}\) onto the additive Fourier modes in \(E_p\).
The raw matrix with entries \(b_p(x-y)\) is \(t_0\Pi_p\).
Furthermore,
\[
 \Pi_p\mathbf1=0,\qquad
 \|\Pi_pF_p-F_p\|\le\eps_0\sqrt p,
 \qquad
 \|\Pi_pF_p\|^2\ge(1-\eps_0^2)p.
\]
The latter assertions follow by Parseval from
\eqref{eq:supply-concentrated-spectrum}.

Let \(B_p^{(2)}\) denote the raw matrix with entries \(b_p(x-y)^2\).
Its additive Fourier eigenvalue at \(s\) is
\[
 \frac{t_0^2}{p}\#\{(h,h')\in E_p^2:h+h'=s\}.
\]
Consequently
\begin{equation}\label{eq:supply-square-kernel}
 \|B_p^{(2)}\|\le t_0^2\eps_0,\qquad
 |b_p(x)|\le t_0\eps_0.
\end{equation}
For later use, if \(e_p^*=|E_p|/p\), direct orthogonality gives
\begin{align}
 \frac1{p-1}\sum_{x\ne0}b_p(x)
     &=-\frac{t_0e_p^*}{p-1},\nonumber\\
 \frac1{p-1}\sum_{x\ne0}b_p(x)^2
     &=\frac{t_0^2(e_p^*-(e_p^*)^2)}{p-1}.
          \label{eq:supply-unit-moments}
\end{align}
In particular both unit means are \(O(\eps_0/p)\).

For a kernel \(h(x-y)\), let \(h\) also denote its raw matrix, and set
\[
 u_S=|S_p|^{-1}\mathbf1_{S_p},\qquad
 u_T=|T_p|^{-1}\mathbf1_{T_p}.
\]
The map from \(\C\oplus\mathcal H_{T_p}\) to
\(\C\oplus\mathcal H_{S_p}\) defined by
\[
 L_p[h]=
 \begin{pmatrix}
 u_S^*hu_T&u_S^*hQ_{T_p}\\
 Q_{S_p}hu_T&Q_{S_p}hQ_{T_p}
 \end{pmatrix}
\]
represents the kernel on the coordinates
\((1,Q_{S_p}\mathbf e_x)\), \((1,Q_{T_p}\mathbf e_y)\).
Indeed \(u_S+Q_{S_p}\mathbf e_x=\mathbf e_x\) for \(x\in S_p\),
and likewise on \(T_p\).
Define
\[
 L_p(u,v)=L_p[(1+ub_p)(1+vb_p)],\qquad
 U_p(u,v)=\frac1{p-1}\sum_{x\ne0}(1+ub_p(x))(1+vb_p(x)).
\]

\begin{lemma}\label{lem:supply-local-contraction}
For sufficiently small \(\eps_0\), uniformly for \(p\in\mathcal S\)
and complex \(u,v\) with \(|u|,|v|\le1.03\),
\[
 \|L_p(u,v)\|\le1+C/p,\qquad |U_p(u,v)|\le1+C/p.
\]
For all sufficiently large such primes,
\begin{equation}\label{eq:src31}
 \|L_p(1,1)\|\le U_p(1,1)(1-1.6/p),
 \qquad U_p(1,1)=1+O(\eps_0/p).
\end{equation}
All constants are absolute.
\end{lemma}

\begin{proof}
The constant kernel 1 has coordinate matrix
\(\begin{psmallmatrix}1&0\\0&0\end{psmallmatrix}\).
In the full field,
\[
 u_S=p^{-1}(\mathbf1+\sqrt{\kappa_p}F_p),\qquad
 u_T=p^{-1}(\mathbf1-F_p/\sqrt{\kappa_p}).
\]
Since \(\Pi_p\) is self-adjoint and kills constants,
\[
 u_S^*\Pi_pu_T=-p^{-2}\|\Pi_pF_p\|^2.
\]
Also \(Q_{S_p}F_p=Q_{T_p}F_p=0\). Hence the two side blocks
arising from \(\Pi_p\) have norms \(O(\eps_0/\sqrt p)\);
here and below balance bounds \(\kappa_p^{\pm1/2}\) absolutely.
Disjoint supports give \(Q_{S_p}Q_{T_p}=0\), and therefore
\[
 \|Q_{S_p}\Pi_pQ_{T_p}\|
 =\|Q_{S_p}(\Pi_p-\tfrac12 I)Q_{T_p}\|\le\tfrac12.
\]
Using \eqref{eq:supply-square-kernel} and
\(\|u_S\|,\|u_T\|=O(p^{-1/2})\), the square kernel contributes
\(O(\eps_0/p)\) to the scalar block,
\(O(\eps_0/\sqrt p)\) to the side blocks, and \(O(\eps_0)\)
to the lower-right block.

Thus, if \(s_p(u,v)\) denotes the scalar block, then
\[
 s_p(u,v)=1+O(1/p),\qquad
 \|\text{lower-right block}\|\le1.03t_0+O(\eps_0)<.95,
\]
uniformly on the stated polydisc, and both side blocks are
\(O(\eps_0/\sqrt p)\).
At \(u=v=1\) the scalar block is real and
\[
 s_p(1,1)
 =1-\frac{2t_0}{p^2}\|\Pi_pF_p\|^2+O(\eps_0/p)
 \le1-\frac{1.7-O(\eps_0)}p.
\]
It is positive for large \(p\).

For completeness, bound the norm of a block matrix by the norm of
the \(2\)-by-\(2\) matrix of its block norms. In the present situation
this is in turn bounded by the largest eigenvalue of
\[
 \begin{pmatrix}
 |s_p(u,v)|&C\eps_0/\sqrt p\\
 C\eps_0/\sqrt p&.95
 \end{pmatrix}.
\]
Because \(|s_p(u,v)|=1+O(1/p)\), its gap from .95 is bounded
below for large \(p\). The eigenvalue formula consequently gives
the upper bound
\[
 |s_p(u,v)|+O(\eps_0^2/p).
\]
This proves the asserted uniform norm estimate and the sharper
bound \(1-(1.7-O(\eps_0))/p\) at \((1,1)\).
Equation~\eqref{eq:supply-unit-moments} gives the estimates for \(U_p\).
After decreasing the fixed \(\eps_0\), the last sharper bound is at
most \(U_p(1,1)(1-1.6/p)\). In particular \(U_p(1,1)>0\).
\end{proof}

\subsection{Tensoring and truncating the kernel}

Choose a sufficiently large fixed \(C_K>0\), as specified below, and put
\(K=\lceil C_KL\rceil\). Define the nonnegative function
\begin{equation}\label{eq:supply-weight}
 W(n)=
 \left(\sum_{\substack{I\subseteq\mathcal S\\|I|\le K}}
              \prod_{p\in I}b_p(n)\right)^2.
\end{equation}
For a polynomial \(P(u,v)\), scalar- or operator-valued, let
\(\mathcal R_K P\) be the sum of its Taylor coefficients with degree
at most \(K\) in each variable, evaluated at \((1,1)\).
Set
\[
 \mathcal L(u,v)=\bigotimes_{p\in\mathcal S}L_p(u,v),\qquad
 \mathcal U(u,v)=\prod_{p\in\mathcal S}U_p(u,v),\qquad
 U=\mathcal U(1,1).
\]
The matrix \(\mathcal R_K\mathcal L\) represents \(W(a-d)\) between
the full tensor coordinates. Similarly,
\(\mathcal R_K\mathcal U\) is the mean of \(W\) on independent
uniform unit residues.

Let \(r_0=1.03\). Lemma~\ref{lem:supply-local-contraction} and
\(H\le .85L+o(L)\) bound the norms of both polynomials on
\(|u|,|v|\le r_0\) by \(\exp(C_0L)\) with an absolute \(C_0\).
Cauchy's coefficient formula applies to the finite-dimensional
operator spaces as well as to scalars: if \(P_{ij}\) is the
coefficient of \(u^iv^j\), then
\(\|P_{ij}\|\le\exp(C_0L)r_0^{-i-j}\).
Summing the geometric series over \(i>K\) or \(j>K\), we obtain
\begin{equation}\label{eq:supply-truncation-error}
 \|\mathcal R_K\mathcal L-\mathcal L(1,1)\|
 +|\mathcal R_K\mathcal U-U|
 \ll \exp(C_0L)r_0^{-K}.
\end{equation}
We also have \(U=\exp(O(\eps_0L))>0\), and
\[
 \|\mathcal L(1,1)\|
 \le U\prod_{p\in\mathcal S}(1-1.6/p)
 \le Ue^{-1.6H}.
\]
Choose the fixed \(C_K\) so large that the right side of
\eqref{eq:supply-truncation-error} is \(o(Ue^{-1.6H})\).
For example, after bounding all the displayed absolute constants,
one can ensure an error at most \(e^{-5L}\).

There is no loss from the dimension of these tensor spaces.
To see explicitly which coordinates are used, write
\[
 V_A=\frac1{m_A}\sum_{a\in A_0}
       \bigotimes_{p\in\mathcal S}(1,Q_{S_p}\mathbf e_{a\bmod p}),
 \qquad
 V_D=\frac1{m_D}\sum_{d\in D_0}
       \bigotimes_{p\in\mathcal S}(1,Q_{T_p}\mathbf e_{d\bmod p}).
\]
The tensor space is the orthogonal direct sum of its centered subset
modes. The squared norm of the mode indexed by the primes of \(t\)
is \(B_A(t)\), respectively \(B_D(t)\).
Each monomial retained in \(\mathcal R_K\mathcal L\) involves
at most \(2K\) primes. At every other prime the local constant kernel
has only a scalar block. Consequently
\(\mathcal R_K\mathcal L=P_A(\mathcal R_K\mathcal L)P_D\),
where \(P_A,P_D\) retain only subset modes with at most \(2K\)
centered factors. Equation~\eqref{eq:src30}, with any fixed
\(C_*>2C_K\), gives
\[
 \|P_AV_A\|^2\le e^{o(L)}R/m_A,\qquad
 \|P_DV_D\|^2\le e^{o(L)}R/m_D.
\]
The kernel identity, the operator norm bound, and Cauchy--Schwarz now
give
\begin{equation}\label{eq:src32}
 \frac1{m_Am_D}\sum_{a\in A_0}\sum_{d\in D_0}W(a-d)
 \le
 U\exp(-1.6H+o(L))
           \left(\frac{R^2}{m_Am_D}\right)^{1/2}.
\end{equation}
Every assertion of this subsection is uniform under the deletion of
one prime from \(\mathcal S\). We may therefore make such a deletion,
and redefine the polynomials and \(W\) accordingly, when addressing
an exceptional character below.

\subsection{The prime mean of the weight}

Set
\begin{equation}\label{eq:supply-character-range}
 \log Q_b=2K\exp(.9L).
\end{equation}
Products of at most \(2K\) primes from \(\mathcal S\) are at most
\(Q_b\). We first record explicitly the analytic estimate used to
average our weight.

\begin{lemma}\label{lem:supply-total-prime-error}
Fix \(C_K>0\), and let \(\phi\) be a fixed smooth function compactly
supported in \((1/2,1)\). Among primitive characters of conductor
at most \(Q_b\), omit the possible exceptional character in the
Landau--Page theorem. Include the conductor-\(1\) principal character.
For every fixed \(D_1>0\),
\begin{equation}\label{eq:src33}
 \sum_{\chi}^{\mathrm{nonexc}}
 \left|
   \sum_{n\ge1}\Lambda(n)\chi(n)\phi(n/X)
       -\ind_{\chi=1}X\int_0^\infty\phi(t)\,dt
 \right|
 \ll_{D_1,\phi,C_K}Xe^{-D_1L}.
\end{equation}
\end{lemma}

\begin{proof}
We give the uniform details because the sum in
\eqref{eq:src33} is unweighted over the primitive characters.
Put \(Q=Q_b\), \(T=Q^5\), and \(M=\log X=e^L\).
The zero-density theorem of Montgomery
\cite[Theorem~1]{Montgomery1969}, in the precise form recorded in
\cite[Lemma~1]{Inoue}, states that
\begin{equation}\label{eq:supply-density-theorem}
 \sum_{\substack{q\le Q\\\chi\bmod q}}^{*}
       N(\sigma,T,\chi)
 \ll
 (Q^2T)^{3(1-\sigma)/(2-\sigma)}
       \bigl(\log(QT)\bigr)^{13}
 \quad
 \left(Q\ge1,\ T\ge2,\ \tfrac12\le\sigma\le1\right).
\end{equation}
Here the star restricts to primitive characters, and \(N\) counts
nontrivial zeros, with multiplicity, having real part at least
\(\sigma\) and ordinate of absolute value at most \(T\).
There is no additional restriction relating \(Q\) and \(T\).

 The standard zero-free region and Landau--Page theorem, in the
 form \cite[Lemma~7.1]{FordGreenKonyaginMaynardTao2018},
 give an absolute \(c>0\)
such that, after omitting at most one primitive real character of
conductor at most \(Q\), every zero in this family with
\(|\Im\rho|\le T\) satisfies
\[
 1-\Re\rho\ge u_0:=\frac{c}{\log(QT)}.
\]
One may choose the omitted character by the Landau--Page theorem
at height zero with parameter \(Q\); decreasing the absolute \(c\)
makes the displayed common strip valid up to height \(T\).
The omitted object is the entire character, not just one zero.

Let \(\mathcal F\) be the retained primitive nonprincipal characters.
For \(u_0\le y\le1/2\), \eqref{eq:supply-density-theorem} implies
\[
 F(y):=\sum_{\chi\in\mathcal F}N(1-y,T,\chi)
 \ll B e^{Ay},\qquad
 B=(\log(QT))^{13},\quad A=3\log(Q^2T).
\]
There are no zeros counted for \(y<u_0\).
Since \(\log Q=2K e^{.9L}=o(M)\), eventually \(M\ge2A\).
Integration by parts for this counting function yields
\begin{align*}
 \sum_{\chi\in\mathcal F}
 \sum_{\substack{|\Im\rho|\le T\\\Re\rho\ge1/2}}
       X^{\Re\rho-1}
 &\le e^{-M/2}F(1/2)
       +M\int_{u_0}^{1/2}e^{-My}F(y)\,dy\\
 &\ll B\exp(-Mu_0/2).
\end{align*}
Now \(\log B=O(L)\), whereas
\[
 Mu_0\gg\frac{e^{.1L}}K.
\]
It follows that this last bound is smaller than \(e^{-DL}\)
for every fixed \(D>0\), for sufficiently large \(L\).

To connect zeros with the required smooth sums, write
\[
 \Phi(s)=\int_0^\infty\phi(t)t^{s-1}\,dt.
\]
It is entire and, uniformly for \(-1/2\le\Re s\le2\),
\(\Phi(s)\ll_{\phi,j}(1+|\Im s|)^{-j}\) for every integer \(j\ge0\).
Mellin inversion, followed by shifting the integral of
\(-L'/L(s,\chi)X^s\Phi(s)\) from \(\Re s=2\) to
\(\Re s=-1/2\), gives for primitive nonprincipal \(\chi\bmod q\)
\begin{equation}\label{eq:supply-smoothed-explicit}
 \sum_n\Lambda(n)\chi(n)\phi(n/X)
 =-\sum_{\rho}X^\rho\Phi(\rho)
      -\ind_{\chi(-1)=1}\Phi(0)
      +O_\phi(X^{-1/2}\log(2q)).
\end{equation}
The zero sum is over nontrivial zeros with multiplicity and
converges absolutely. For clarity, the residue at zero is the
simple trivial zero of an even nonprincipal character.
On \(\Re s=-1/2\), the functional equation expresses the
logarithmic derivative through its absolutely convergent counterpart
on \(\Re s=3/2\), together with gamma factors. It gives the uniform
bound \(O(\log(q(2+|\Im s|)))\) on this line.
The decay of \(\Phi\) proves the displayed remainder.
Shifting first at suitable finite heights and then passing to the
limit justifies the contour operation using the ordinary zero-count
bound \(N(0,t,\chi)\ll (t+1)\log(q(t+2))\);
these standard functional-equation and explicit-formula facts are
given in \cite[Chapter~10]{MontgomeryVaughan2007}; see also
\cite[Sections~3.7.1 and~3.7.3]{Helfgott}.
In particular the error contains no constant depending on the
distance of a possible exceptional zero from 1.

We sum absolute values in \eqref{eq:supply-smoothed-explicit}.
The contribution from zeros with \(\Re\rho\ge1/2\) and
\(|\Im\rho|\le T\) is
\[
 O_\phi(XB e^{-Mu_0/2}).
\]
There are \(O(Q^2T\log(QT))\) zeros in the retained family at these
heights, so those with real part below \(1/2\) contribute at most
\(O_\phi(X^{1/2}Q^2T\log(QT))\).
Using \(|\Phi(\beta+i\gamma)|\ll_\phi(1+|\gamma|)^{-3}\)
and summing the zero-count bound in dyadic height intervals, the
zeros above \(T\) contribute
\(O_\phi(XQ^2T^{-2}\log(QT))\).
Finally the remainders in \eqref{eq:supply-smoothed-explicit} sum to
\[
 O_\phi\bigl(Q^2(1+X^{-1/2}\log(2Q))\bigr).
\]

Each of these terms is \(O_{D_1,\phi,C_K}(Xe^{-D_1L})\).
For the first, use the already proved estimate on \(Mu_0\).
For the second and fourth, use \(Q^7=X^{o(1)}\).
For the third, \(Q^2T^{-2}=Q^{-8}\), whose logarithm is a
negative quantity of order \(K e^{.9L}\).
The conductor-\(1\) character is added by the classical prime
number theorem with its zero-free-region error, followed by smooth
partial summation; its error is
\(O_\phi(Xe^{-c\sqrt{\log X}})\), which is also sufficient.
This proves \eqref{eq:src33}.
\end{proof}

We now perform the possible deletion promised after
\eqref{eq:src32}, using the one-prime removal of an exceptional conductor
as in \cite[Section~7, following Lemma~7.1]{FordGreenKonyaginMaynardTao2018}.
If the exceptional primitive character of
conductor at most \(Q_b\) has conductor equal to a product of primes
in \(\mathcal S\), delete one of those primes from \(\mathcal S\).
Otherwise make no deletion. Keep \(K,Q_b\) fixed and redefine
\(W,\mathcal L,\mathcal U,U,H\) using the resulting set.
Equations~\eqref{eq:supply-sparse-mass} and \eqref{eq:src32}
continue to hold. Every primitive character induced by a monomial
of \(W\) has squarefree conductor formed from primes of
\(\mathcal S\); hence none of those induced characters is the
exceptional one.

Fix from now on a nonzero nonnegative smooth \(\phi\) compactly
supported in \((1/2,1)\).
For \(j=1,2\) and a multiplicative character \(\lambda\bmod p\),
write
\[
 c_{p,j}(\lambda)
 =\frac1{p-1}\sum_{x\ne0}b_p(x)^j\overline{\lambda(x)}.
\]
Multiplicative Fourier inversion expresses \(b_p(x)^j\) as
\(\sum_\lambda c_{p,j}(\lambda)\lambda(x)\) on \(\F_p^*\).
Equations~\eqref{eq:supply-square-kernel} and
\eqref{eq:supply-unit-moments}, followed by Cauchy--Schwarz, give
\begin{equation}\label{eq:supply-mellin-coefficients}
 |c_{p,j}(\lambda)|\ll p^{-1/2}\quad(\lambda\ne1,\ j=1,2),
 \qquad |c_{p,j}(1)|\ll p^{-1}\quad(j=1,2).
\end{equation}
For example the mean of \(b_p^2\) on units is \(O(1/p)\), and
the mean of \(b_p^4\) is at most
\(\|b_p\|_\infty^2\) times that mean, also \(O(1/p)\).
All constants here are absolute.

Expand \eqref{eq:supply-weight} by its two subsets \(I,J\), and
expand each participating local factor multiplicatively.
For an integer coprime to every prime of \(\mathcal S\), group
the resulting terms by the primitive character they induce.
This gives
\begin{equation}\label{eq:supply-grouped-expansion}
 W(n)=\sum_{\substack{\chi\ {\rm primitive}\\
                         \operatorname{cond}(\chi)\le Q_b}}
             C(\chi)\chi(n)
 \qquad
 \left((n,\prod_{p\in\mathcal S}p)=1\right).
\end{equation}
Only squarefree conductors supported on \(\mathcal S\) occur.
The conductor bound follows because \(|I\cup J|\le2K\).

For a fixed primitive character \(\chi\) of conductor \(d\),
its local nonprincipal character is prescribed at every \(p\mid d\).
There are three possible statuses for such a prime: membership
in \(I\) alone, in \(J\) alone, or in both.
The sum of the absolute coefficient costs at that prime is
at most \(C/\sqrt p\) by \eqref{eq:supply-mellin-coefficients}.
For \(p\nmid d\), the absent status contributes 1 and the other
three statuses use principal coefficients, with total cost
at most \(1+C/p\).
Discarding the degree restrictions only enlarges this nonnegative
majorant. Therefore even the sum of the absolute values of all
terms grouped into \(C(\chi)\) is bounded by
\begin{equation}\label{eq:supply-grouped-bound}
 \prod_{p\mid d}\frac C{\sqrt p}
       \prod_{\substack{p\in\mathcal S\\p\nmid d}}(1+C/p)
 \le e^{C'L}.
\end{equation}
In the last inequality all conductor-prime factors are at most 1
for sufficiently large \(L\), and \(\sum_{p\in\mathcal S}1/p=O(L)\).
This is a uniform bound for each grouped coefficient; we will
multiply it by the total error in \eqref{eq:src33}.
The principal coefficient retains its signs and satisfies exactly
\[
 C(1)=\mathcal R_K\mathcal U
      =U+o(Ue^{-1.6H})=(1+o(1))U.
\]

Every prime in the support of \(\phi(p/X)\) is larger than every
prime in \(\mathcal S\), because
\(\exp(\exp(.9L))=X^{o(1)}\). Thus
\eqref{eq:supply-grouped-expansion} holds on all desired prime
inputs. Apply \eqref{eq:src33} to the grouped sum, taking its fixed
\(D_1\) larger than the absolute constants in
\eqref{eq:supply-grouped-bound} and in \(U=\exp(O(\eps_0L))\).
No exceptional character occurs, by our deletion.
The resulting error is \(o(UX)\).

To pass from von Mangoldt sums to primes, it is harmless that the
grouped identity was only asserted on units. Indeed the number
of primitive characters of conductor at most \(Q_b\) is at most
\(Q_b^2\), and
\[
 \sum_{\substack{n\le X\\n\ {\rm a\ prime\ power}\\
                              n\ {\rm not\ prime}}}\Lambda(n)
 \ll X^{1/2}(\log X)^2.
\]
The absolute contribution of all these terms to the grouped
character sum is at most
\[
 e^{C'L}Q_b^2 X^{1/2}(\log X)^2
       =X^{1/2+o(1)}=o(UX).
\]
This bound also covers prime powers whose prime base lies in
\(\mathcal S\), for which extension from unit classes can change
the value of the grouped expansion. We have proved
\begin{equation}\label{eq:src34}
 \sum_{p\ {\rm prime}}(\log p)\phi(p/X)W(p)
   =(1+o(1))UX\int_0^\infty\phi(t)\,dt.
\end{equation}

\subsection{Using coverage of every sufficiently large prime}

The elementary bound \(|b_p(n)|\le1\) gives, with
\(M_{\mathcal S}=|\mathcal S|\),
\[
 W(n)\le
  (K+1)^2(1+M_{\mathcal S})^{2K}
  \le\exp\bigl(O(K\exp(.9L))\bigr)=X^{o(1)}
\]
uniformly in \(n\).
By eventual prime coverage, every prime counted in
\eqref{eq:src34}, for sufficiently large \(L\), has at least one
representation \(p=a-d\) with \(a\in A\), \(d\in D\).
Since \(a,-d\ge0\) and \(p<X\), every such representation has
\(a,-d\le X\).

Remove all primes possessing a representation with
\(a\le R+N_*\) or \(-d\le R+N_*\).
The number of removed primes is at most
\[
 A(R+N_*)B(X)+A(X)B(R+N_*)=X^{3/4+o(1)}
\]
by \eqref{eq:src1}. Their contribution to the left side of
\eqref{eq:src34} is still \(X^{3/4+o(1)}\), by the uniform
bound on \(W\), boundedness of \(\phi\), and \(\log p\le\log X\).
This is \(o(UX)\).
Every remaining prime has a representation by a pair in
\(A_0\times D_0\). Choose one such pair for each prime.
Different primes give different pairs, and the other pairs have
nonnegative weight. Consequently, for a fixed \(c_\phi>0\),
\begin{equation}\label{eq:supply-coverage-lower}
 \frac1{m_Am_D}\sum_{a\in A_0,d\in D_0}W(a-d)
 \ge
 c_\phi U\frac{X}{m_Am_D\log X}.
\end{equation}
Here we used
\((\log p)\phi(p/X)\le\|\phi\|_\infty\log X\)
to remove the prime weight from the surviving lower bound in
\eqref{eq:src34}. This step uses coverage of the full prime set.

Put \(z=X/(m_Am_D)=R^2/(m_Am_D)\).
Equation~\eqref{eq:supply-tail-size} gives \(z\ge e^{-o(L)}\).
Comparing \eqref{eq:supply-coverage-lower} with
\eqref{eq:src32}, and canceling the positive \(U\), gives
\[
 c_\phi z e^{-L}
 \le e^{-1.6H+o(L)}\sqrt z.
\]
Since \(H\ge .68L\), this implies
\[
 \sqrt z\le\exp\bigl(-.088L+o(L)\bigr),
\]
contradicting \(z\ge e^{-o(L)}\).
Thus \eqref{eq:src28} cannot fail along an unbounded sequence,
and Proposition~\ref{prop:supply} follows.

\begin{remark}[An optional zero-free refinement]
Theorem~1.1 of \cite{OpenAIQRH2026} implies that
\eqref{eq:src33} also holds when the sum includes \emph{all} primitive
characters of conductor at most \(Q_b\), including the conductor-\(1\)
principal character. Indeed, every nontrivial zero then has real part at
most \(7/8\). The smoothed explicit formula
\eqref{eq:supply-smoothed-explicit}, the rapid decay of \(\Phi\), and the
zero-count bound used above give an error
\(O_\phi(X^{7/8}\log(2q))\) for each primitive character of conductor
\(q\); for the conductor-\(1\) character, subtract the pole main term.
There are \(O(Q_b^2)\) such characters, so their total error is
\(O_\phi(Q_b^2X^{7/8}\log(2Q_b))\), and
\[
 Q_b^2X^{7/8}\log(2Q_b)=X^{7/8+o(1)}
   \ll_{D_1,C_K}Xe^{-D_1L}
\]
for every fixed \(D_1>0\), since
\(K=\lceil C_KL\rceil\) and \(\log Q_b=o(\log X)\).
The proof of Proposition~\ref{prop:supply}, including its
exceptional-character deletion, uses the classical estimate of
Lemma~\ref{lem:supply-total-prime-error}, not this zero-free theorem.
\end{remark}
